\input{./._relpath-to-latexroot.ltx}
\documentclass[\latexroot/\projectname]{subfiles}
\input{\latexroot/@resources/subfile-setup.ltx}

\begin{document}

\section{Empirical Analysis}
\whenintegrated{\label{sec:empirical}}

We bring the cost-minimizing allocation model of Section~\ref{sec:model}
to a panel of 32 advanced economies over 1989--2022. The empirical
strategy is a two-stage estimator: a first stage recovers each
country's structural sector size $\hat\delta_{ij}$, and a second stage
uses these as interaction variables to estimate the slope and
threshold behaviour of debt allocation.

%-------------------------------------------------------------
\subsection{Data and sector-size variables}
%-------------------------------------------------------------

\paragraph{Sources.}
\begin{itemize}
  \item \textbf{Debt holdings:} \citet{ArslanalpTsuda2014}, sovereign
        debt by creditor sector.
  \item \textbf{Sector-size proxies:} World Bank,
        \emph{Global Financial Development Database} (GFDD).
  \item \textbf{Panel:} 32 advanced economies, 1989--2022, for a
        total of $N=780$ country--year observations. Coverage
        thickens after 2000; the full 32-country panel is available
        from 2007 onward.
\end{itemize}

\paragraph{Country--year proxies for structural sector size $S_j$.}
\begin{align*}
s_{h,it}  &= \tfrac{\text{Pension fund assets}_{it}+\text{Insurance assets}_{it}}{\text{GDP}_{it}}
            && \to\; S_h \text{ (non-bank capacity)} \\
s_{pb,it} &= \tfrac{\text{Bank deposit funding}_{it}}{\text{GDP}_{it}}
            && \to\; S_{pb} \text{ (private-bank size)} \\
s_{cb,it} &= \tfrac{\text{Central bank assets}_{it}}{\text{GDP}_{it}}
            && \to\; S_{cb} \text{ (central-bank size)}
\end{align*}

\paragraph{Foreign sector.}
No country-specific size proxy is available---foreign creditors are
global balance sheets, not a local sector we can size in GFDD. We
therefore run \emph{no first-stage regression for foreigners}. Foreign
debt is treated as the accounting residual
\begin{equation}\label{eq:foreign-residual}
  d_{f,it}=D_{it}-d_{h,it}-d_{cb,it}-d_{pb,it}.
\end{equation}

%-------------------------------------------------------------
\subsection{Empirical strategy: two-stage estimator}
%-------------------------------------------------------------

\paragraph{Threshold.}
$\bar D = 100\%$ of GDP, chosen from the data: banking's marginal
absorption accelerates sharply once $D_{it}$ crosses this level
(discussed in the Motivation).

\paragraph{Stage 1---structural sector size $\hat\delta_{ij}$.}
For each country $i$ and sector $j\in\{h,cb,pb\}$, regress sector
share on total debt with country and year fixed effects, allowing a
kink at $\bar D$:
\begin{equation}\label{eq:stage1}
s_{ijt}=\delta_{ij}+\delta_{jt}
       +\alpha_{1j}D_{it}
       +\alpha_{2j}\,\mathbf{1}\{D_{it}\!\ge\!\bar D\}\,(D_{it}-\bar D)
       +u_{ijt}.
\end{equation}
We read $\hat\delta_{ij}$, the country fixed effect, as the part of
country $i$'s sector-$j$ size that does not move with time, after
stripping out variation explained by $D_{it}$ and the kink.

\paragraph{Stage 2---debt allocation.}
The effective slope is a product of a baseline rate and a size
scaling, with a kink at $\bar D$. In its original factored form,
\begin{equation}\label{eq:stage2-factored}
d^j_{it}
= \underbrace{\beta_j\!\left(1+\gamma_j\hat\delta_{ij}\right)}_{\text{baseline slope}} D_{it}
+ \mathbf{1}_{it}\,
  \underbrace{\eta_j\!\left(1+\rho_j\hat\delta_{ij}\right)}_{\text{post-threshold slope change}}
  (D_{it}-\bar D)
+ \varepsilon_{ijt}.
\end{equation}
Expanding the two products gives the estimating equation reported in
the tables:
\begin{equation}\label{eq:stage2}
d^j_{it}=\beta_j D_{it}
        +\theta_j(\hat\delta_{ij}\!\cdot\! D_{it})
        +\mu_j\mathbf{1}_{it}(D_{it}-\bar D)
        +\nu_j\mathbf{1}_{it}(\hat\delta_{ij}\!\cdot\!(D_{it}-\bar D))
        +\varepsilon_{ijt},
\end{equation}
with $\theta_j=\beta_j\gamma_j$, $\mu_j=\eta_j$, and
$\nu_j=\eta_j\rho_j$.
\begin{itemize}
  \item $\beta_j$: baseline slope---sector $j$'s absorption per unit
        of $D$.
  \item $\theta_j$: \emph{size-scaling}---larger $\hat\delta_{ij}$
        implies steeper absorption.
  \item $\mu_j,\nu_j$: post-threshold slope change and its
        size-scaling.
\end{itemize}

%-------------------------------------------------------------
\subsection{First-stage results: structural sector sizes}
%-------------------------------------------------------------

Table~\ref{tab:stage1} reports estimates of \eqref{eq:stage1}.

\begin{table}[htbp]\centering
\caption{First stage: structural sector sizes}
\label{tab:stage1}
\begin{tabular}{lccc}
\toprule
 & (1) & (2) & (3) \\
 & Non-bank & Central bank & Private bank \\
 & $s_{iht}$ & $s_{icbt}$ & $s_{ipbt}$ \\
\midrule
$\alpha_{1j}$: $D_{it}$
  & $-0.203$ & $0.016$ & $-0.321$ \\
  & $(0.234)$ & $(0.044)$ & $(0.265)$ \\
$\alpha_{2j}$: highD $\!\cdot\!(D_{it}-\bar D)$
  & $-0.184$ & $0.520^{***}$ & $0.856^{**}$ \\
  & $(0.203)$ & $(0.167)$ & $(0.372)$ \\
\midrule
$N$           & 547 & 718 & 732 \\
Groups        & 31 & 32 & 32 \\
$R^2$ (within)& 0.436 & 0.603 & 0.199 \\
Year / Country FE & Yes & Yes & Yes \\
\bottomrule
\multicolumn{4}{p{0.8\linewidth}}{\footnotesize Notes: Standard errors in parentheses. *** $p<0.01$, ** $p<0.05$, * $p<0.10$.}
\end{tabular}
\end{table}

\paragraph{Takeaway.}
\begin{itemize}
  \item \textbf{Non-banks.} $\alpha_{1h},\alpha_{2h}$ are
        insignificant: non-bank size does not move with total debt,
        consistent with a \emph{hard} capacity $\bar\omega S_h$.
  \item \textbf{Central- and private-bank size.} $\alpha_{2j}>0$ and
        significant: banking-sector size \emph{grows} once
        $D_{it}>\bar D$, consistent with banking's no-hard-ceiling
        property in the model.
\end{itemize}

%-------------------------------------------------------------
\subsection{Second-stage results: debt allocation}
%-------------------------------------------------------------

Table~\ref{tab:stage2} reports estimates of \eqref{eq:stage2}.

\begin{table}[htbp]\centering
\caption{Second stage: debt allocation by sector}
\label{tab:stage2}
\setlength{\tabcolsep}{8pt}
\renewcommand{\arraystretch}{1.15}
\begin{tabular}{lccc}
\toprule
 & Non-bank ($d_h$) & Central bank ($d_{cb}$) & Private bank ($d_{pb}$) \\
\midrule
$\beta_j$: $D_{it}$
  & $0.314^{***}$ & $0.053^{***}$ & $0.192^{***}$ \\
  & $(0.040)$ & $(0.012)$ & $(0.020)$ \\
$\theta_j$: $\hat{\delta}_{ij}\!\cdot D_{it}$
  & $0.0018^{***}$ & $0.0131^{***}$ & $-0.0002$ \\
  & $(0.0005)$ & $(0.0044)$ & $(0.0007)$ \\
$\mu_j$: threshold slope
  & $-0.018$ & $0.334^{**}$ & $0.462^{**}$ \\
  & $(0.171)$ & $(0.153)$ & $(0.220)$ \\
$\nu_j$: threshold $\times\,\hat{\delta}$
  & $-0.004$ & $-0.055$ & $-0.009$ \\
  & $(0.003)$ & $(0.036)$ & $(0.008)$ \\
\midrule
$N$  & $753$ & $780$ & $780$ \\
$F$  & $83.26$ & $177.55$ & $72.93$ \\
$R^2$& $0.802$ & $0.705$ & $0.819$ \\
\bottomrule
\multicolumn{4}{p{0.85\linewidth}}{\footnotesize Notes: Standard errors in parentheses. *** $p<0.01$, ** $p<0.05$, * $p<0.10$.}
\end{tabular}
\end{table}

\paragraph{Takeaways.}
\begin{itemize}
  \item \textbf{All sectors absorb part of new debt.} $\beta_j>0$
        across the three sectors.
  \item \textbf{Non-banks scale with size.} $\theta_h>0$ and
        significant: countries with intrinsically larger non-bank
        sectors absorb a larger share of government debt at any given
        $D$. This maps to model prediction~2 (size scaling).
  \item \textbf{Banking takes over above $\bar D$.}
        $\mu_{cb}+\mu_{pb}\approx 0.80$: banking's marginal
        absorption jumps by roughly 80 percentage points once
        $D_{it}>100\%$ of GDP. This maps to model prediction~4
        (banking takeover).
\end{itemize}

\paragraph{Foreign debt residual.}
Using the fitted domestic components from \eqref{eq:stage2}, the
implied foreign residual tracks the data: correlation $=0.779$, slope
$=0.744^{***}$, $R^2=0.608$, MAE $=8.79$.

%-------------------------------------------------------------
\subsection{Concavity in debt absorption}
%-------------------------------------------------------------

\paragraph{Specification.}
For sector $j\in\{h,f\}$,
\begin{equation}\label{eq:concavity}
d^j_{it}=\beta_j D_{it}+\kappa_j D_{it}^2
        +\theta_j(\hat\delta_h\!\cdot D_{it})+\varepsilon_{it},
\end{equation}
where $\theta_j$ is included only for the non-bank sector. Pooled and
country-FE estimates are reported separately, with and without
$\hat\delta_h\!\cdot D$.

Table~\ref{tab:concavity} reports estimates of \eqref{eq:concavity}.

\begin{table}[htbp]\centering
\caption{Concavity in debt absorption}
\label{tab:concavity}
\setlength{\tabcolsep}{6pt}
\renewcommand{\arraystretch}{1.15}
\begin{tabular}{lcccccc}
\toprule
& \multicolumn{4}{c}{\textbf{Non-bank} ($d_h$)}
& \multicolumn{2}{c}{\textbf{Foreign} ($d_f$)} \\
\cmidrule(lr){2-5}\cmidrule(lr){6-7}
& \multicolumn{2}{c}{Spec C: no $\hat\delta_h$}
& \multicolumn{2}{c}{Spec B: with $\hat\delta_h\!\cdot D$}
& \multicolumn{2}{c}{No $\hat\delta_h\!\cdot D$} \\
\cmidrule(lr){2-3}\cmidrule(lr){4-5}\cmidrule(lr){6-7}
& Pooled & FE & Pooled & FE & Pooled & FE \\
\midrule
$\beta\,(D)$
& $0.343^{***}$ & $0.402^{***}$ & $0.359^{***}$ & $0.299^{*}$
& $0.387^{***}$ & $0.392^{***}$ \\
$\theta\,(\hat\delta_h\!\cdot D)$
& -- & -- & $0.0022^{***}$ & $0.0011^{*}$
& -- & -- \\
$\kappa\,(D^2)$
& $-0.0003$ & $-0.0016^{***}$ & $-0.0010$ & $-0.0009$
& $-0.0014^{***}$ & $-0.0009$ \\
\midrule
$N$   & $780$ & $780$ & $701$ & $701$ & $747$ & $747$ \\
$R^2$ & $0.735$ & $0.381$ & $0.801$ & $0.381$ & $0.773$ & $0.454$ \\
\bottomrule
\multicolumn{7}{p{0.95\linewidth}}{\footnotesize Notes: Standard errors omitted for space. *** $p<0.01$, ** $p<0.05$, * $p<0.10$.}
\end{tabular}
\end{table}

\paragraph{Takeaways.}
\begin{itemize}
  \item \textbf{Non-banks.} Concavity ($\kappa_h<0$) shows up
        \emph{within country} but disappears once we control for
        size ($\hat\delta_h\!\cdot D$): the slowdown is driven by
        country-specific capacity $S_h$.
  \item \textbf{Foreigners.} Concavity ($\kappa_f<0$) shows up
        \emph{across countries} but not within: consistent with a
        \emph{common} foreign ceiling $\bar D_f$.
\end{itemize}

\paragraph{Map to predictions.}
These two findings jointly support model prediction~3: the marginal
absorption shares of non-banks and foreigners fall as $D$ grows, but
the source of concavity differs by sector---country-specific $S_h$
for non-banks, a common $\bar D_f$ for foreigners.

% Smart bibliography: only included when compiled standalone with citations
\smartbib

\end{document}
